\section{低次元の場合と次数付き同型}
\label{sec:lowdim}

大域次元が3または4の場合には、最小生成元の次数だけから
定理\ref{thm:ore}に必要なAS indexを決定できる。
以下では両側 noetherian性から得られる成長度の下界と、最小自由分解の双対性を用いる。

\subsection{最小自由分解とHilbert級数の零点}

$A$を大域次元$d$、AS index$\ell$の連結次数付きAS-regular代数とし、
$k_A=A/A_{\geq 1}$の最小次数付き右自由分解を
\[
  0\longrightarrow P_d\longrightarrow\cdots\longrightarrow P_0
  \longrightarrow k_A\longrightarrow0,
  \qquad
  P_i=\bigoplus_{\nu=1}^{b_i}A(-a_{i\nu})
\]
と書く。ここで$A(-a)_m=A_{m-a}$である。
AS双対性によれば、$P_0=A$、$P_d=A(-\ell)$であり、
$P_{d-i}$のシフトの多重集合は
$\{\ell-a_{i\nu}\}_{\nu=1}^{b_i}$に一致する
\cite[Lemma 1.1]{LPWZ}。
また
\begin{equation}
  a_{i\nu}\geq i
  \label{eq:lowdim-shift-bound}
\end{equation}
である。実際、最小性より$P_i\otimes_A k\cong\Tor_i^A(k,k)$であり、
$\Tor_i^A(k,k)$を計算する被約bar複体の第$i$項
$A_{\geq1}^{\otimes i}$は内部次数$i$未満の成分を持たない。

自由分解の各次数でEuler標数を取ると
\begin{equation}
  H_A(t)=\frac{1}{p_A(t)},
  \qquad
  p_A(t)=\sum_{i=0}^{d}(-1)^i\sum_{\nu=1}^{b_i}t^{a_{i\nu}}
  \label{eq:lowdim-euler}
\end{equation}
を得る。以下で用いる代数は次数1と2の元で生成される。
この場合、$t=1$における$p_A(t)$の零点の位数を$m$とすると
$\GKdim A\leq m$である。ここでは$p_A(1)\ne0$のとき$m=0$とする。
この不等式は、$h_j=\dim_k A_j$と置いたときの
\[
  \sum_{j=0}^{N}h_j
  \leq eH_A(e^{-1/N})=O(N^m)
\]
から従う。AS-regular性に含まれる有限GK次元の条件により、
この式のHilbert級数は$0<t<1$で収束する。
また、$1$と固定した最小生成元が張る空間を$U$とすれば
$U^N\subseteq\bigoplus_{j=0}^{2N}A_j$なので、
上の評価はGK次元の定義に適用できる。
両側 noetherianかつ$d\geq3$ならば
\cite[Lemma 1.2]{LPWZ}より$\GKdim A\geq3$であり、したがって
\begin{equation}
  p_A(1)=p_A'(1)=p_A''(1)=0
  \label{eq:lowdim-three-zeros}
\end{equation}
となる。

\subsection{生成元の次数が\texorpdfstring{$(1,1,1,2)$}{(1,1,1,2)}の場合}

\begin{proposition}
\label{prop:index-four}
$A$を両側 noetherianな連結次数付きAS-regular$\C$代数とする。
$\gldim A=4$であり、最小生成元の次数が$(1,1,1,2)$ならば、
\[
  \ASind A=5,
  \qquad
  H_A(t)=\frac{1}{(1-t)^3(1-t^2)}.
\]
さらに、$B=\C\langle A_1\rangle\subseteq A$は
3次元のquadratic Koszul AS-regular代数であり、ある
$\sigma\in\Aut_{\mathrm{gr}}(B)$と次数2の$\sigma$導分$\delta$に対して
\[
  A\cong B[w;\sigma,\delta],\qquad \deg w=2
\]
となる。この同型は$B$上の恒等写像を延長する。
\end{proposition}

\begin{proof}
AS双対性で最小分解の両端を決め、
\eqref{eq:lowdim-three-zeros}から中間項のシフトの平方和を求める。
各シフトが入る区間上の不等式がAS indexの上界を与える。

AS indexを$\ell$と書く。最小生成元の次数とAS双対性より
\[
\begin{split}
  P_1&=A(-1)^{\oplus3}\oplus A(-2),\\
  P_3&=A(-\ell+1)^{\oplus3}\oplus A(-\ell+2),
  \qquad P_4=A(-\ell).
\end{split}
\]
$P_3$にはシフト$\ell-2$の直和因子があるので、
\eqref{eq:lowdim-shift-bound}から$\ell\geq5$を得る。
$P_2=\bigoplus_{j=1}^{r}A(-s_j)$と書けば、同じ評価と双対性により
\begin{equation}
  2\leq s_j\leq\ell-2,
  \qquad
  \{s_1,\ldots,s_r\}=\{\ell-s_1,\ldots,\ell-s_r\}
  \label{eq:four-middle-shifts}
\end{equation}
が多重集合として成立する。
\eqref{eq:lowdim-euler}の分母は
\[
  p_A(t)=1-3t-t^2+\sum_{j=1}^{r}t^{s_j}
       -3t^{\ell-1}-t^{\ell-2}+t^\ell.
\]
$p_A(1)=0$から$r=6$であり、
\eqref{eq:four-middle-shifts}の対称性から
\begin{equation}
  \sum_{j=1}^{6}s_j=3\ell
  \label{eq:four-first-moment}
\end{equation}
を得る。さらに$p_A''(1)=0$は
\[
  -2+\sum_{j=1}^{6}s_j(s_j-1)
  -3(\ell-1)(\ell-2)-(\ell-2)(\ell-3)+\ell(\ell-1)=0
\]
であるから、\eqref{eq:four-first-moment}を代入して
\begin{equation}
  \sum_{j=1}^{6}s_j^2=3\ell^2-10\ell+14
  \label{eq:four-second-moment}
\end{equation}
を得る。一方、$2\leq s_j\leq\ell-2$より
\[
  s_j(\ell-s_j)-2(\ell-2)
   =(s_j-2)(\ell-2-s_j)\geq0.
\]
これを$j$について加え、
\eqref{eq:four-first-moment}と\eqref{eq:four-second-moment}を使うと
\[
  12\ell-24\leq
  \sum_{j=1}^{6}s_j(\ell-s_j)=10\ell-14.
\]
したがって$\ell\leq5$であり、先の下界と合わせて$\ell=5$となる。
このとき$s_j\in\{2,3\}$かつ$\sum_js_j=15$なので、
次数2と次数3の最小関係式はそれぞれ3本である。
よって
\[
  p_A(t)=1-3t+2t^2+2t^3-3t^4+t^5
        =(1-t)^3(1-t^2).
\]

定理\ref{thm:ore}を$n=3$に適用すると、
$B=\C\langle A_1\rangle$は3次元のquadratic Koszul AS-regular代数であり、
主張のOre表示が得られる。
\end{proof}

\subsection{生成元の次数が\texorpdfstring{$(1,1,2)$}{(1,1,2)}の場合}

\begin{proposition}
\label{prop:index-three}
$A$を両側 noetherianな連結次数付きAS-regular$\C$代数とする。
$\gldim A=3$であり、最小生成元の次数が$(1,1,2)$ならば、
\[
  \ASind A=4,
  \qquad
  H_A(t)=\frac{1}{(1-t)^2(1-t^2)}.
\]
さらに$B=\C\langle A_1\rangle$は、次数1の生成元$x,y$を適切に選ぶと
\begin{equation}
  B_q=\C\langle x,y\rangle/(yx-qxy)\quad(q\in\C^\times),
  \qquad\text{または}\qquad
  B_J=\C\langle x,y\rangle/(yx-xy-x^2)
  \label{eq:plane-normal-forms}
\end{equation}
のいずれかである\cite[\S 2, Corollary 2.2の後]{RZ}。
ある$\sigma\in\Aut_{\mathrm{gr}}(B)$と次数2の$\sigma$導分$\delta$に対して、
$B$上の恒等写像を延長する次数付き同型
\[
  A\cong B[w;\sigma,\delta],\qquad\deg w=2
\]
が存在する。
\end{proposition}

\begin{proof}
AS indexを$\ell$と書く。最小分解の双対性から
\[
  p_A(t)=1-2t-t^2+2t^{\ell-1}+t^{\ell-2}-t^\ell
\]
である。\eqref{eq:lowdim-three-zeros}のうち$p_A'(1)=0$を用いると
\[
  0=-2-2+2(\ell-1)+(\ell-2)-\ell=2\ell-8
\]
となるので、$\ell=4$である。したがって
\[
  p_A(t)=1-2t+2t^3-t^4=(1-t)^2(1-t^2).
\]
定理\ref{thm:ore}を$n=2$に適用すると、
$B=\C\langle A_1\rangle$は2次元AS-regular代数であり、
主張のOre表示と$H_B(t)=(1-t)^{-2}$が得られる。

ここで$B$の表示を確認する。
$B$は次数1の2元$x,y$で最小生成されるので、
AS双対性と$H_B(t)=(1-t)^{-2}$から最小右自由分解は
\[
  0\longrightarrow B(-2)\xrightarrow{d_2}B(-1)^{\oplus2}
   \xrightarrow{d_1}B\longrightarrow k_B\longrightarrow0
\]
となる。$d_1(e_1)=x$、$d_1(e_2)=y$とし、
$d_2(1)=e_1a+e_2b$と書くと、$a,b\in B_1$であり、
$B$の唯一の最小関係式は
$r=x\otimes a+y\otimes b\in B_1\otimes B_1$である。
この分解に$\operatorname{Hom}_B(-,B)$を施すと、AS条件から
\[
  B/(Ba+Bb)\cong k
\]
が得られる。従って$a,b$は$B_1$の基底であり、
$r$の係数行列は可逆である。

$r$の可換化を$\overline r\in\operatorname{Sym}^2(B_1)$とする。
$\overline r=0$なら$r$は$xy-yx$の非零スカラー倍であり、
$B\cong\C[x,y]$、すなわち$q=1$の量子平面となる。
$\overline r$の階数が2ならば、$\C$上の変数変換で
$\overline r$を$xy$の非零スカラー倍にできる。
この座標では$r=\alpha xy+\beta yx$となり、
係数行列の可逆性により$\alpha\beta\ne0$だから、
$q=-\alpha/\beta$と置けば量子平面の表示を得る。
$\overline r$の階数が1なら、スカラー倍と変数変換により
\[
  r=x^2+\gamma(xy-yx)
\]
と書ける。係数行列が可逆なので$\gamma\ne0$であり、
$y'=\gamma y$と置くと$y'x-xy'-x^2=0$となり、Jordan平面の表示を得る。
これで\eqref{eq:plane-normal-forms}の全場合を得た。
\end{proof}

\begin{corollary}
\label{cor:lowdim-converse}
$n=2$の場合は$B$を\eqref{eq:plane-normal-forms}のいずれかとし、
$n=3$の場合は$B$を次数1の3元で生成される両側 noetherianな
quadratic Koszul AS-regular代数で$\gldim B=3$を満たすものとする。
$\sigma\in\Aut_{\mathrm{gr}}(B)$と次数2の$\sigma$導分$\delta$を任意に取る。
このとき$B[w;\sigma,\delta]$は両側 noetherianなAS-regular代数であり、
その大域次元は$n+1$、最小生成元の次数は$(1^n,2)$である。
\end{corollary}

\begin{proof}
量子平面とJordan平面は、それぞれ
$\C[x][y;\tau]$、$\C[x][y;\partial]$というOre拡大である。
ここで$\tau(x)=qx$、$\partial(x)=x^2$であり、$\deg x=\deg y=1$とする。
これらは両側 noetherianな2次元AS-regular代数であり、Hilbert級数は$(1-t)^{-2}$である。
これらの性質と、一般の$A=B[w;\sigma,\delta]$のAS-regular性には、
次のOre拡大定理を適用できる。
noetherian AS-regularな連結次数付き代数はtwisted Calabi--Yauであり
\cite[\S 1.2]{LWW}、自己同型によるOre拡大はその次元を1増やす
\cite[Theorem 3.3]{LWW}。
得られた代数のAS条件と有限大域次元は\cite[Lemma 1.3]{LWW}から従い、
両側 noetherian性はOre拡大に対するHilbertの基底定理による
\cite[\S 1.3]{LWW}。
これらを$\C[x]$からの拡大にも適用すれば、最初に述べた2つの平面の性質も得られる。

いずれの$n$についても$H_B(t)=(1-t)^{-n}$である。
$n=3$のときは、Koszul性とAS双対性による線形分解
$0\to B(-3)\to B(-2)^3\to B(-1)^3\to B\to k\to0$から従う。
$A$の左$B$自由基底が$\{w^j\}_{j\geq0}$なので
\[
  H_A(t)=\frac{H_B(t)}{1-t^2}
        =\frac{1}{(1-t)^n(1-t^2)},
\]
従って有限GK次元の条件も満たされる。
$A_1=B_1$、$A_2=B_2\oplus\C w$であるから、
$B$の最小生成元に$w$を加えたものは最小生成系である。
\end{proof}

\subsection{Oreデータの同値関係}

次数1で生成される部分代数$\C\langle A_1\rangle$は
次数付き同型によって保たれる。
したがって、上述のOre表示の同型判定は次の形に尽きる。

\begin{proposition}
\label{prop:ore-isomorphism}
$B,B'$を次数1で生成される連結次数付き$\C$代数とし、
\[
  A=B[w;\sigma,\delta],\qquad
  A'=B'[w';\sigma',\delta'],\qquad
  \deg w=\deg w'=2
\]
を次数付きOre拡大とする。
ここで$\sigma,\sigma'$は次数付き自己同型、
$\delta,\delta'$は次数2の対応する歪導分である。
次数付き$\C$代数同型$F:A\to A'$は、
次のデータと1対1に対応する：
次数付き同型$\phi:B\to B'$、$\lambda\in\C^\times$、$u\in B'_2$であって
\begin{align}
  \phi\circ\sigma&=\sigma'\circ\phi,
  \label{eq:ore-iso-sigma}\\
  \phi\bigl(\delta(b)\bigr)
    &=\lambda\delta'\bigl(\phi(b)\bigr)
      +u\phi(b)-\sigma'\bigl(\phi(b)\bigr)u
      \qquad(b\in B)
  \label{eq:ore-iso-delta}
\end{align}
を満たすものである。
この対応は
\[
  F|_B=\phi,\qquad F(w)=\lambda w'+u
\]
で与えられる。
\end{proposition}

\begin{proof}
$A_1=B_1$、$A'_1=B'_1$なので、次数付き同型$F$は
$\C\langle A_1\rangle=B$を$B'$へ写し、次数付き同型$\phi:B\to B'$を誘導する。
$A_2=B_2\oplus\C w$、$A'_2=B'_2\oplus\C w'$であるから、
ある一意な$\lambda\in\C^\times$と$u\in B'_2$により
$F(w)=\lambda w'+u$と書ける。
$wb=\sigma(b)w+\delta(b)$に$F$を施し、左$B'$自由基底の$w'$係数と定数項を比較すると
\begin{align*}
  (\lambda w'+u)\phi(b)
   &=\lambda\sigma'\bigl(\phi(b)\bigr)w'
     +\lambda\delta'\bigl(\phi(b)\bigr)+u\phi(b),\\
  \phi\bigl(\sigma(b)\bigr)(\lambda w'+u)+\phi\bigl(\delta(b)\bigr)
   &=\lambda\phi\bigl(\sigma(b)\bigr)w'
     +\phi\bigl(\sigma(b)\bigr)u+\phi\bigl(\delta(b)\bigr)
\end{align*}
より\eqref{eq:ore-iso-sigma}と\eqref{eq:ore-iso-delta}を得る。

逆にこれらの等式が成り立てば、$F|_B=\phi$、$F(w)=\lambda w'+u$は
Ore関係式を保つので、次数付き準同型$F:A\to A'$を定める。
その像は$B'$および$w'=\lambda^{-1}(F(w)-u)$を含むため、$F$は全射である。
また、$b_m\ne0$である元$\sum_{i=0}^{m}b_iw^i$の像を
左$B'$自由基底で展開したとき、$(w')^m$の係数は
$\lambda^m\phi(b_m)\ne0$である。
従って$F$は単射でもある。
\end{proof}

\eqref{eq:ore-iso-sigma}と\eqref{eq:ore-iso-delta}は$B_1$上で確認すれば十分である。
前者の両辺は代数準同型であり、前者が成立すると、後者の両辺は
$B$から$B'$への同じ$\phi\circ\sigma$および$\phi$に関する歪導分になるためである。
具体的には、その差を$D$とすると
$D(ab)=\phi(\sigma(a))D(b)+D(a)\phi(b)$であり、
$D|_{B_1}=0$から積の長さに関する帰納法で$D=0$を得る。
特に同じ$B$の上で変数を$\widetilde w=\lambda w+u$（$u\in B_2$）に取り替えると、
$\sigma$は変わらず、歪導分は
\[
  \widetilde\delta(b)=\lambda\delta(b)+ub-\sigma(b)u
  \qquad(b\in B)
\]
に変わる。
